\section{Trace-sealed primitive recursion}
\label{sec:trace}

Adaptive algorithms and relational constructions are the point at which fixed
proof graphs meet variable object parameters.  A proof cannot contain one
existential-elimination rule for each of an object-variable number \(N\) of
steps.  SCR packages the entire interaction as one finite trace certificate
and reflects it by one quantifier-free induction.

\subsection{Normalized trace data}

Fix a query \(q\), horizon \(N\), action bound \(M(q,N)\), and an action code
\(\tau<M^N\).  Let \(a_i=\mathsf{digit}_M(\tau,i)\).  Define by primitive
recursion
\[
\begin{aligned}
  \State(q,\tau,0)&=\mathsf{init}(q),\\
  \State(q,\tau,i+1)&=\mathsf{next}(q,i,\State(q,\tau,i),a_i),\\
  \Bad(q,\tau,0)&=0,\\
  \Bad(q,\tau,i+1)&=\Bad(q,\tau,i)+
     \mathsf{fail}(q,i,\State(q,\tau,i),a_i).
\end{aligned}
\]
The failure bit includes the digit bound, the local relational checker, and
any resource side condition.  The trace validator is the quantifier-free
equation
\[
       \Trace(q,N,\tau)\;:\!\Longleftrightarrow\;
       \Bad(q,\tau,N)=0.
\]
A final-state predicate may be included in \(\mathsf{fail}\) at a terminal
stage or checked separately.

\subsection{Prefix reflection}

Let \(I(q,s)\) be a reflected invariant bit.  Assume object derivations of:
\[
\begin{aligned}
  &I(q,\mathsf{init}(q))=1,\\
  &I(q,s)=1\wedge\mathsf{fail}(q,i,s,a)=0
      \;\longrightarrow\;
    I(q,\mathsf{next}(q,i,s,a))=1,\\
  &\Bad(q,\tau,i+1)=0
      \;\longrightarrow\;
    \Bad(q,\tau,i)=0\wedge\mathsf{fail}(q,i,\State_i,a_i)=0.
\end{aligned}
\]
The third line is a reusable arithmetic lemma about a sum of nonnegative
Boolean values.

\begin{theorem}[Trace reflection]
\label{thm:trace}
iPRA proves
\[
  \Trace(q,N,\tau)\;\longrightarrow\;
  I(q,\State(q,\tau,N))=1.
\]
\end{theorem}

\begin{proof}
Apply quantifier-free induction to
\[
  \Phi(i)\equiv
  \Bad(q,\tau,i)=0\to I(q,\State(q,\tau,i))=1.
\]
At zero use the initialization theorem.  At a successor, assume
\(\Bad(i+1)=0\); split it into \(\Bad(i)=0\) and local failure zero, invoke the
induction hypothesis, then the local preservation theorem.  Every displayed
formula is a Boolean combination of equalities and inequalities between
\(\PTm\) terms, so \(\Phi.\qf\) is obtained structurally.  Instantiate the
result at \(N\).
\end{proof}

\subsection{Sealed trace theorem}

If a producer proves \(\exists\tau\,\Trace(q,N,\tau)\), open \(\tau\), apply
\cref{thm:trace}, derive a final goal not mentioning \(\tau\), and seal it.
No trace code escapes.  If the final result includes the computed endpoint,
two alternatives are available:
\begin{enumerate}[label=(\alph*)]
\item expose the relational output
\(\exists\tau\,[\Trace(q,N,\tau)\wedge o=O(\State_N)]\);
\item prove observer independence or select a bounded least trace and define a
      term-valued endpoint.
\end{enumerate}

\subsection{Changing precision and state-dependent demand}

The state record may contain the current precision, mesh, error budget,
remaining fuel, and any finite summary of previous observations.  Therefore a
changing precision does not change the induction formula; it changes a field
of \(\State_i\).  The fixed invariant is a relation between those fields.  This
resolves the common packet mistake of trying to perform a separate induction
whose formula changes with the stage.

A reverse demand pass can itself be normalized as a trace.  Start from the
consumer's final demand, apply primitive-recursive transfer functions
backwards through a fixed graph, and store the resulting resource vector in the
initial state of the forward replay.  If demands are cyclic, use inflationary
iteration on a bounded product lattice.

\begin{proposition}[Bounded cyclic demand closure]
Suppose \(F:[0,R_1]\times\cdots\times[0,R_k]\to
[0,R_1]\times\cdots\times[0,R_k]\) is componentwise inflationary.  Starting
at \(0\) and iterating \(F\), a fixed point is reached in at most
\(1+\sum_j R_j\) stages.
\end{proposition}

\begin{proof}
If a stage is not fixed, at least one component increases by at least one, so
the rank \(\sum_j x_j\) strictly increases.  It is bounded by
\(\sum_jR_j\).  All comparisons and the rank invariant are quantifier-free,
so the proof is a qf-induction template in iPRA.
\end{proof}

\subsection{Producer modes for traces}

There are three distinct trace producers.
\begin{description}
\item[Deterministic.]  The action is a PR function of the state.  No external
certificate is needed; the state recursion itself is the trace.
\item[Boundedly searchable.]  A complete trace is below an explicit cutoff and
has a reflected verifier.  Use whole-trace least search if a name is required,
or seal its existence otherwise.
\item[Relational boundary.]  The application proves or assumes only
\(\exists\tau\,\Trace\), perhaps from a deep finite-cover or descent theorem.
Keep this exact existential as the producer boundary.
\end{description}
